\subsection{Electron Energy Resolution}\label{sec:ELECTRON_ENERGY_RESOLUTION}
As a figure of merit for the process, we can compute the energy resolution from the cheated calibrated energy of the main electron. This is typically the value that analyzers compute when estimating the energy resolution of a \ac{TPC}. It is obtained for an idealized clustering with ad hoc calorimetry and a high hit threshold (in this case we use 3 for a typical track). The energy resolution is computed as the \ac{RMS} of the relative error of the reconstructed variable versus the truth. The results are shown in \cref{fig:RMS}. The formula for the energy resolution fit is:

%residuals = np.sqrt(np.power(p2, 2) + np.power(p1 / np.sqrt(x - b), 2) + np.power(p0 / (x - b), 2))
\begin{equation}
    \Omega_{(E_\text{true})} = \sqrt{\left(\frac{p_2}{\sqrt{E_{\text{true}} - b}}\right)^2 + \left(\frac{p_1}{E_{\text{true}} - b}\right)^2 + p_0^2},
\end{equation}
\noindent where \( p_0 \), \( p_1 \), and \( p_2 \) are the parameters of the fit, and \( b \) is a constant offset and corresponds to the detection threshold. These parameters can be interpreted as the intrinsic resolution evaluated to \num{11(1)}\%, the stochastic term, and the noise term, respectively.

\begin{figure}[ht!]
    \centering
    \includegraphics[width=\linewidth]{Images/hd_1x2x6_centralAPA_marley_Cheated_Energy_Resolution_NHits3.png}
    \caption[width=1\linewidth]{Energy resolution from calibrated charge with cheated clustering and hit threshold of 3.\label{fig:RMS}}
\end{figure}

Applying this calibration to the main electron cluster, we obtain the final calibrated energy of the main electron. For all the adjacent clusters, we simply apply the charge correction factor obtained from the main electron cluster, assuming that these clusters are gamma-induced ``blips'' and therefore have a similar response to the main electron. \cref{fig:SIGNAL_DAUGHTER_SMEARING} shows the smearing matrices of the main electron and gamma energies after the calibration. Additionally, the smearing matrix of the visible energy (defined as the sum of the main electron and gamma energy) is shown.

\begin{figure}[ht!]
    \centering
    \includegraphics[width=\linewidth]{Images/hd_1x2x6_centralAPA_marley_True&Reco_Daughter_Energy.png}
    \caption[width=1\linewidth]{Signal daughter smearing obtained after main electron calibration.\label{fig:SIGNAL_DAUGHTER_SMEARING}}
\end{figure}

\textbf{Electron energy resolution with and without flash-matching.}